% GENERATED FILE. Edit canonical lessons, then run npm run generate:books.
% canonical-source-hash: fnv1a64:fad40a0ba7a77731
% canonical-lessons: AR-C503-mutasamih, AR-C503-asabi, AR-C503-marih, AR-C503-munazzam, AR-C503-muhmil, AR-R503-review-496-503

\chapter{Tempers and habits}
\label{ch:ar-tempers-and-habits}
\hlchaptermodality{503}

\begin{chapteropening}
\textbf{By the end of this chapter:} \emph{I can describe a person as tolerant, irritable, cheerful, organized or careless.}

Complete the last lesson of chapter 503: I can describe a person as tolerant, irritable, cheerful, organized or careless.
\end{chapteropening}

\section[\texorpdfstring{\ar{متسامح} (mutasāmiḥ)}{متسامح (mutasāmiḥ)}]{\ar{متسامح} (mutasāmiḥ) — tolerant, forgiving}

\label{lesson:AR-C503-mutasamih}



\begin{quote}
Before the new one: say the Arabic for affectionate, then the Arabic for friendly.
\end{quote}

\subsection*{You'll want to know: \ar{متسامح}}
\textbf{\ar{متسامح}} — \emph{mutasāmiḥ} — \textquotedblleft{}tolerant, forgiving\textquotedblright{}.

\begin{quote}
\ar{أبي} \ar{رجل} \ar{متسامح} \ar{ولا} \ar{يغضب} \ar{كثيرا}.
\end{quote}

\begin{quote}
\emph{abī rajul mutasāmiḥ wa-lā yaghḍabu kathīran.}
\end{quote}

\begin{quote}
My father is a tolerant man and does not get angry often.
\end{quote}

\begin{cousinweb}
\textbf{\ar{متسامح}} (\emph{mutasāmiḥ}) is built on the native Arabic root \emph{s-m-ḥ}, \textquotedblleft{}to be generous, to allow\textquotedblright{} (Lane). Samaḥa, to allow, and sāmaḥa, to forgive, which you know, share the root.
\end{cousinweb}

\begin{grammarlens}[title={What you've built}]
One more word for describing what you lived and felt.
\end{grammarlens}

\subsection*{Your turn}
\begin{itemize}
\raggedright
  \item \emph{Say it:} \emph{mutasāmiḥ}
  \item \emph{Say it:} \emph{mutasāmiḥ}, once more
  \item \emph{Say it:} \emph{abī rajul mutasāmiḥ wa-lā yaghḍabu kathīran.} — My father is a tolerant man and does not get angry often.
  \item \emph{Recall:} say the Arabic for affectionate, then the Arabic for friendly, then say \emph{mutasāmiḥ} again
\end{itemize}

\begin{tcolorbox}[breakable,colback=teal!4,colframe=teal!35!black,title={Before you move on}]
What does \ar{متسامح} mean? (\textquotedblleft{}Tolerant, forgiving\textquotedblright{}.) Say it once more.
\end{tcolorbox}
\section[\texorpdfstring{\ar{عصبي} (ʿaṣabī)}{عصبي (ʿaṣabī)}]{\ar{عصبي} (ʿaṣabī) — irritable, edgy}

\label{lesson:AR-C503-asabi}



\begin{quote}
Before the new one: say the Arabic for friendly, then the Arabic for tolerant.
\end{quote}

\subsection*{You'll want to know: \ar{عصبي}}
\textbf{\ar{عصبي}} — \emph{ʿaṣabī} — \textquotedblleft{}irritable, edgy\textquotedblright{}.

\begin{quote}
\ar{يصبح} \ar{المدير} \ar{عصبيا} \ar{قبل} \ar{الاجتماع}.
\end{quote}

\begin{quote}
\emph{yuṣbiḥu al-mudīr ʿaṣabiyyan qabla al-ijtimāʿ.}
\end{quote}

\begin{quote}
The manager gets edgy before the meeting.
\end{quote}

It is built on ʿaṣab, a nerve: a nervy person.

\begin{cousinweb}
\textbf{\ar{عصبي}} (\emph{ʿaṣabī}) is built on the native Arabic root \emph{ʿ-ṣ-b}, \textquotedblleft{}to bind, to tie tight\textquotedblright{} (Lane). ʿAṣabī is the adjective of ʿaṣab, a nerve.
\end{cousinweb}

\begin{grammarlens}[title={What you've built}]
One more word for describing what you lived and felt.
\end{grammarlens}

\subsection*{Your turn}
\begin{itemize}
\raggedright
  \item \emph{Say it:} \emph{ʿaṣabī}
  \item \emph{Say it:} \emph{ʿaṣabī}, once more
  \item \emph{Say it:} \emph{yuṣbiḥu al-mudīr ʿaṣabiyyan qabla al-ijtimāʿ.} — The manager gets edgy before the meeting.
  \item \emph{Recall:} say the Arabic for friendly, then the Arabic for tolerant, then say \emph{ʿaṣabī} again
\end{itemize}

\begin{tcolorbox}[breakable,colback=teal!4,colframe=teal!35!black,title={Before you move on}]
What does \ar{عصبي} mean? (\textquotedblleft{}Irritable, edgy\textquotedblright{}.) Say it once more.
\end{tcolorbox}
\section[\texorpdfstring{\ar{مرح} (mariḥ)}{مرح (mariḥ)}]{\ar{مرح} (mariḥ) — cheerful, fun-loving}

\label{lesson:AR-C503-marih}



\begin{quote}
Before the new one: say the Arabic for tolerant, then the Arabic for irritable.
\end{quote}

\subsection*{You'll want to know: \ar{مرح}}
\textbf{\ar{مرح}} — \emph{mariḥ} — \textquotedblleft{}cheerful, fun-loving\textquotedblright{}.

\begin{quote}
\ar{خالي} \ar{رجل} \ar{مرح} \ar{ويضحك} \ar{كثيرا}.
\end{quote}

\begin{quote}
\emph{khālī rajul mariḥ wa-yaḍḥaku kathīran.}
\end{quote}

\begin{quote}
My uncle is a cheerful man and laughs a lot.
\end{quote}

The noun \ar{مرح}, read maraḥ, is merriment.

\begin{cousinweb}
\textbf{\ar{مرح}} (\emph{mariḥ}) is built on the native Arabic root \emph{m-r-ḥ}, \textquotedblleft{}to be lively, merry\textquotedblright{} (Lane). Mariḥ is the adjective; maraḥ, with the same letters, is the noun.
\end{cousinweb}

\begin{grammarlens}[title={What you've built}]
One more word for describing what you lived and felt.
\end{grammarlens}

\subsection*{Your turn}
\begin{itemize}
\raggedright
  \item \emph{Say it:} \emph{mariḥ}
  \item \emph{Say it:} \emph{mariḥ}, once more
  \item \emph{Say it:} \emph{khālī rajul mariḥ wa-yaḍḥaku kathīran.} — My uncle is a cheerful man and laughs a lot.
  \item \emph{Recall:} say the Arabic for tolerant, then the Arabic for irritable, then say \emph{mariḥ} again
\end{itemize}

\begin{tcolorbox}[breakable,colback=teal!4,colframe=teal!35!black,title={Before you move on}]
What does \ar{مرح} mean? (\textquotedblleft{}Cheerful, fun-loving\textquotedblright{}.) Say it once more.
\end{tcolorbox}
\section[\texorpdfstring{\ar{منظم} (munaẓẓam)}{منظم (munaẓẓam)}]{\ar{منظم} (munaẓẓam) — organized, tidy}

\label{lesson:AR-C503-munazzam}



\begin{quote}
Before the new one: say the Arabic for irritable, then the Arabic for cheerful.
\end{quote}

\subsection*{You'll want to know: \ar{منظم}}
\textbf{\ar{منظم}} — \emph{munaẓẓam} — \textquotedblleft{}organized, tidy\textquotedblright{}.

\begin{quote}
\ar{مكتب} \ar{أختي} \ar{منظم} \ar{ونظيف}.
\end{quote}

\begin{quote}
\emph{maktab ukhtī munaẓẓam wa-naẓīf.}
\end{quote}

\begin{quote}
My sister's desk is tidy and clean.
\end{quote}

\begin{cousinweb}
\textbf{\ar{منظم}} (\emph{munaẓẓam}) is built on the native Arabic root \emph{n-ẓ-m}, \textquotedblleft{}to string pearls in order\textquotedblright{} (Lane). Munaẓẓam is the passive participle of Form II: set in order.
\end{cousinweb}

\begin{grammarlens}[title={What you've built}]
One more word for describing what you lived and felt.
\end{grammarlens}

\subsection*{Your turn}
\begin{itemize}
\raggedright
  \item \emph{Say it:} \emph{munaẓẓam}
  \item \emph{Say it:} \emph{munaẓẓam}, once more
  \item \emph{Say it:} \emph{maktab ukhtī munaẓẓam wa-naẓīf.} — My sister's desk is tidy and clean.
  \item \emph{Recall:} say the Arabic for irritable, then the Arabic for cheerful, then say \emph{munaẓẓam} again
\end{itemize}

\begin{tcolorbox}[breakable,colback=teal!4,colframe=teal!35!black,title={Before you move on}]
What does \ar{منظم} mean? (\textquotedblleft{}Organized, tidy\textquotedblright{}.) Say it once more.
\end{tcolorbox}
\section[\texorpdfstring{\ar{مهمل} (muhmil)}{مهمل (muhmil)}]{\ar{مهمل} (muhmil) — careless}

\label{lesson:AR-C503-muhmil}



\begin{quote}
Before the new one: say the Arabic for cheerful, then the Arabic for organized.
\end{quote}

\subsection*{You'll want to know: \ar{مهمل}}
\textbf{\ar{مهمل}} — \emph{muhmil} — \textquotedblleft{}careless\textquotedblright{}.

\begin{quote}
\ar{كسر} \ar{العامل} \ar{المهمل} \ar{الزجاج}.
\end{quote}

\begin{quote}
\emph{kasara al-ʿāmil al-muhmil az-zujāj.}
\end{quote}

\begin{quote}
The careless worker broke the glass.
\end{quote}

\begin{cousinweb}
\textbf{\ar{مهمل}} (\emph{muhmil}) is built on the native Arabic root \emph{h-m-l}, \textquotedblleft{}to leave neglected\textquotedblright{} (Lane). Muhmil is the participle of Form IV: one who leaves things neglected.
\end{cousinweb}

\begin{grammarlens}[title={What you've built}]
One more word for describing what you lived and felt.
\end{grammarlens}

\subsection*{Your turn}
\begin{itemize}
\raggedright
  \item \emph{Say it:} \emph{muhmil}
  \item \emph{Say it:} \emph{muhmil}, once more
  \item \emph{Say it:} \emph{kasara al-ʿāmil al-muhmil az-zujāj.} — The careless worker broke the glass.
  \item \emph{Recall:} say the Arabic for cheerful, then the Arabic for organized, then say \emph{muhmil} again
\end{itemize}

\begin{tcolorbox}[breakable,colback=teal!4,colframe=teal!35!black,title={Before you move on}]
What does \ar{مهمل} mean? (\textquotedblleft{}Careless\textquotedblright{}.) Say it once more.
\end{tcolorbox}
\section[(dialogue)]{Words from chapters 496–503 — a review}

\label{lesson:AR-R503-review-496-503}



\begin{quote}
Say the Arabic for excited and to get excited.
\end{quote}

\subsection*{Your turn}
Say these aloud:

\begin{itemize}
\raggedright
  \item excited, to get excited, relief, at ease, to feel relieved
  \item tense, to get flustered, agitated, hesitant, to hold back
  \item to make happy, to sadden, to make laugh, to make angry, to blame
  \item honest, upright, untruthful, loyal, faithful
  \item conceited, arrogant, selfish, cowardly, bold
  \item wise, skilful, creative, sensible, confidence
  \item talkative, reserved, nosy, affectionate, friendly
  \item tolerant, irritable, cheerful, organized, careless
\end{itemize}

\begin{tcolorbox}[breakable,colback=teal!4,colframe=teal!35!black,title={Before you move on}]
Excited? (\textbf{\ar{متحمس}}, \emph{mutaḥammis}.) To get excited? (\textbf{\ar{تحمس}}, \emph{taḥammasa}.) Relief? (\textbf{\ar{ارتياح}}, \emph{irtiyāḥ}.) At ease? (\textbf{\ar{مرتاح}}, \emph{murtāḥ}.) To feel relieved? (\textbf{\ar{ارتاح}}, \emph{irtāḥa}.) Tense? (\textbf{\ar{متوتر}}, \emph{mutawattir}.) To get flustered? (\textbf{\ar{ارتبك}}, \emph{irtabaka}.) Agitated? (\textbf{\ar{مضطرب}}, \emph{muḍṭarib}.) Hesitant? (\textbf{\ar{متردد}}, \emph{mutaraddid}.) To hold back? (\textbf{\ar{كتم}}, \emph{katama}.) To make happy? (\textbf{\ar{أسعد}}, \emph{asʿada}.) To sadden? (\textbf{\ar{أحزن}}, \emph{aḥzana}.) To make laugh? (\textbf{\ar{أضحك}}, \emph{aḍḥaka}.) To make angry? (\textbf{\ar{أغضب}}, \emph{aghḍaba}.) To blame? (\textbf{\ar{لام}}, \emph{lāma}.) Honest? (\textbf{\ar{أمين}}, \emph{amīn}.) Upright? (\textbf{\ar{نزيه}}, \emph{nazīh}.) Untruthful? (\textbf{\ar{كاذب}}, \emph{kādhib}.) Loyal? (\textbf{\ar{مخلص}}, \emph{mukhliṣ}.) Faithful? (\textbf{\ar{وفي}}, \emph{wafī}.) Conceited? (\textbf{\ar{مغرور}}, \emph{maghrūr}.) Arrogant? (\textbf{\ar{متكبر}}, \emph{mutakabbir}.) Selfish? (\textbf{\ar{أناني}}, \emph{anānī}.) Cowardly? (\textbf{\ar{جبان}}, \emph{jabān}.) Bold? (\textbf{\ar{جريء}}, \emph{jarīʾ}.) Wise? (\textbf{\ar{حكيم}}, \emph{ḥakīm}.) Skilful? (\textbf{\ar{ماهر}}, \emph{māhir}.) Creative? (\textbf{\ar{مبدع}}, \emph{mubdiʿ}.) Sensible? (\textbf{\ar{عاقل}}, \emph{ʿāqil}.) Confidence? (\textbf{\ar{ثقة}}, \emph{thiqa}.) Talkative? (\textbf{\ar{ثرثار}}, \emph{tharthār}.) Reserved? (\textbf{\ar{متحفظ}}, \emph{mutaḥaffiẓ}.) Nosy? (\textbf{\ar{فضولي}}, \emph{fuḍūlī}.) Affectionate? (\textbf{\ar{حنون}}, \emph{ḥanūn}.) Friendly? (\textbf{\ar{ودود}}, \emph{wadūd}.) Tolerant? (\textbf{\ar{متسامح}}, \emph{mutasāmiḥ}.) Irritable? (\textbf{\ar{عصبي}}, \emph{ʿaṣabī}.) Cheerful? (\textbf{\ar{مرح}}, \emph{mariḥ}.) Organized? (\textbf{\ar{منظم}}, \emph{munaẓẓam}.) Careless? (\textbf{\ar{مهمل}}, \emph{muhmil}.)
\end{tcolorbox}
